\chapter{Cryptography for Systems Programmers}
\label{ch:crypto}

\epigraph{``Cryptography is not a source of security. Cryptography is a tool that amplifies security when used correctly, and amplifies disaster when misused.''}{}

\learninggoal{
	\item Distinguish between cryptographic primitives and understand their correct usage.
	\item Identify common low-level pitfalls in crypto code: side channels, RNG failures, padding oracles.
	\item Understand constant-time programming and why it matters.
	\item Explain how memory-safety bugs in crypto libraries subvert cryptographic guarantees.
	\item Apply the cryptographic engineer's mindset to systems programming.
}

\section{Cryptographic Primitives for Systems Programming}

A systems programmer encounters three core primitive types:

\subsection{Symmetric Encryption}

A single key encrypts and decrypts data.

\[
	c = E(k, m) \quad\quad m = D(k, c)
\]

Common algorithms: AES (in various modes: AES-GCM, AES-CBC, AES-CTR), ChaCha20.

\textbf{Key low-level details:}
\begin{itemize}
	\item AES operates on 16-byte blocks. A single-block encryption is 10--14 rounds of SubBytes, ShiftRows, MixColumns, AddRoundKey. Each round is a sequence of table lookups and XOR operations.
	\item AES-NI instructions (\texttt{aesenc}, \texttt{aesdec}, \texttt{aeskeygenassist}) perform these rounds in hardware, eliminating table-lookup side channels.
	\item Block cipher modes matter enormously: AES-ECB leaks plaintext patterns; AES-CBC requires unpredictable IVs; AES-GCM provides authenticated encryption.
\end{itemize}

\subsection{Asymmetric (Public-Key) Encryption}

Two keys: public for encryption, private for decryption.

\[
	c = E(pk, m) \quad\quad m = D(sk, c)
\]

Common algorithms: RSA, ECDH, Ed25519.

\textbf{Key low-level details:}
\begin{itemize}
	\item RSA operates on integers modulo a large composite \(N = pq\). Encryption is \(c = m^e \bmod N\). Decryption is \(m = c^d \bmod N\). The exponentiation is implemented with square-and-multiply, which is vulnerable to timing side channels if not implemented in constant time.
	\item Elliptic curve operations are point additions and scalar multiplications on elliptic curves over finite fields. The security depends on the hardness of the discrete logarithm problem. Incorrect implementations can compute the wrong curve or fail to validate points.
\end{itemize}

\subsection{Hash Functions}

A deterministic, one-way function with collision resistance.

\[
	h = H(m)
\]

Common algorithms: SHA-256, SHA-3, BLAKE2.

\textbf{Key low-level details:}
\begin{itemize}
	\item SHA-256 processes 64-byte blocks through 64 rounds of compression. The state is 8 32-bit words. Like AES, implementations that use table lookups for the compression function may leak state through cache timing.
	\item Length-extension attacks: SHA-256 (and all Merkle-Damgård hashes) allow computing \(H(m \,\|\, \text{pad}(m) \,\|\, \text{extension})\) given only \(H(m)\) and \(\text{len}(m)\), without knowing \(m\). This is why HMAC exists.
\end{itemize}

\section{Common Low-Level Crypto Pitfalls}

\subsection{Constant-Time Programming}

The most important concept in cryptographic engineering: cryptographic code must take the same amount of time regardless of its inputs.

\begin{definition}[Constant-Time Code]
	\textbf{Constant-time} code executes in a time that depends only on the length of the inputs, not on their values. Any data-dependent timing variation can be used as a side channel to leak secret information.
\end{definition}

Consider a naive string comparison:

\begin{lstlisting}[language=C, caption=Variable-time comparison, label=lst:var-time]
int memcmp_vartime(const void *a, const void *b, size_t n) {
    const unsigned char *pa = a, *pb = b;
    for (size_t i = 0; i < n; i++) {
        if (pa[i] != pb[i])
            return 0;  // Returns early on first mismatch!
    }
    return 1;
}
\end{lstlisting}

This function returns as soon as it finds a mismatch. An attacker who can measure the response time can determine, byte by byte, where the correct key differs from their guess. This is a \emph{timing oracle}.

The constant-time version:

\begin{lstlisting}[language=C, caption=Constant-time comparison]
int memcmp_consttime(const void *a, const void *b, size_t n) {
    const unsigned char *pa = a, *pb = b;
    unsigned char result = 0;
    for (size_t i = 0; i < n; i++) {
        result |= pa[i] ^ pb[i];  // Always processes all bytes
    }
    return (result == 0);
}
\end{lstlisting}

\begin{principle}[No Secret-Dependent Branches or Lookups]
	To write constant-time code: (1) never branch based on secret data, (2) never index memory with secret data (cache side channel), (3) always process all elements.
\end{principle}

\subsection{Random Number Generation}

Cryptography requires unpredictable random numbers. The standard C \texttt{rand()} is not suitable. The Linux kernel exposes random bytes through \texttt{/dev/urandom} (or the \texttt{getrandom} syscall).

\begin{lstlisting}[language=C, caption=Reading random bytes on Linux]
#include <sys/random.h>

void generate_key(uint8_t *key, size_t len) {
    ssize_t n = getrandom(key, len, 0);
    if (n != (ssize_t)len) {
        // Handle error (should not happen with GRND_NONBLOCK)
    }
}
\end{lstlisting}

Common RNG mistakes:
\begin{itemize}
	\item Seeding with \texttt{time(NULL)}: if the attacker knows when the seed was generated, they can reproduce the entire random sequence.
	\item Using \texttt{rand()} or \texttt{random()}: these are pseudo-random and not cryptographically secure.
	\item Reading from \texttt{/dev/random} instead of \texttt{/dev/urandom}: \texttt{/dev/random} blocks when the kernel entropy pool is low, causing denial-of-service.
	\item Failing to check the return value of read/getrandom.
\end{itemize}

\subsection{Padding Oracle Attacks}

A padding oracle attack exploits the fact that decryption failures reveal whether padding was valid:

\begin{lstlisting}[language=C, caption=AES-CBC padding oracle]
int decrypt_and_verify(const uint8_t *ciphertext, size_t len) {
    uint8_t plaintext[MAX_LEN];
    int pt_len = aes_cbc_decrypt(ciphertext, len, key, iv, plaintext);
    if (pt_len < 0)
        return -1;  // Bad padding: attacker learns this!
    if (!verify_mac(plaintext, pt_len))
        return -2;  // Bad MAC
    return pt_len;
}
\end{lstlisting}

An attacker who can distinguish ``bad padding'' from ``bad MAC'' can decrypt arbitrary ciphertexts byte by byte. The fix: always validate MAC first, and return a single error code regardless of what failed.

\section{Memory Safety in Cryptographic Code}

Cryptographic code is particularly sensitive to memory-safety bugs because the security guarantees collapse entirely if an attacker can read or write arbitrary memory.

\subsection{Secret Management}

\begin{itemize}
	\item \textbf{Secrets on the heap:} a secret stored on the heap may survive in freed memory, then be reused by a different allocation.
	\item \textbf{Secrets in core dumps:} a core dump captures all process memory, including private keys.
	\item \textbf{Secrets in swap:} swapped-out pages containing secrets are written to disk.
	\item \textbf{Secrets in side channels:} cache timing, branch prediction, and power analysis can leak secret-dependent access patterns.
\end{itemize}

\begin{lstlisting}[language=C, caption=Secure memory zeroing]
void secure_zero(void *buf, size_t len) {
    volatile unsigned char *p = buf;
    for (size_t i = 0; i < len; i++)
        p[i] = 0;
    // `volatile` prevents the compiler from optimising away the zeroing
}
\end{lstlisting}

Using \texttt{memset(buf, 0, len)} is not secure for clearing secrets: the compiler may optimise it away since the buffer is not read after being written. The \texttt{volatile} qualifier forces the zeroing to be emitted.

\subsection{Library-Level Vulnerabilities}

The most devastating crypto vulnerabilities often come from memory-safety bugs in the library itself rather than cryptographic weaknesses:

\begin{itemize}
	\item \textbf{Heartbleed (CVE-2014-0160):} a missing bounds check in OpenSSL's heartbeat extension allowed reading up to 64 KB of arbitrary process memory---including private keys.
	\item \textbf{CCA\_Decrypt\_OpenSSL:} a use-after-free in OpenSSL's X.509 certificate parsing could lead to code execution when processing a malicious certificate.
	\item \textbf{bleichenbacher oracle:} an implementation bug in RSA PKCS\#1 v1.5 decryption (not a bounds check but a timing variation) allowed breaking RSA encryption.
\end{itemize}

\begin{principle}[Crypto Libraries Are Not Immune to Memory Bugs]
	A cryptographic algorithm may be mathematically unbreakable, but its implementation in C is subject to the same memory-safety bugs as any other C program. The security of a cryptosystem is the security of its weakest component.
\end{principle}

\section{Authenticated Encryption}

The modern approach to symmetric encryption is \emph{authenticated encryption}: encryption that guarantees both confidentiality and integrity. AES-GCM and ChaCha20-Poly1305 are the dominant AEAD (Authenticated Encryption with Associated Data) constructions.

\[
	(c, t) = AEAD\_Enc(k, nonce, ad, plaintext)
\]

Where \(t\) is an authentication tag. On decryption, the tag is verified before the plaintext is released:

\begin{lstlisting}[language=C, caption=AES-GCM decryption must verify tag]
int aead_decrypt(const uint8_t *ciphertext, size_t len,
                 const uint8_t *tag, const uint8_t *key,
                 const uint8_t *nonce, const uint8_t *ad,
                 size_t ad_len, uint8_t *plaintext) {
    // 1. Decrypt and verify tag atomically
    if (crypto_aead_decrypt(plaintext, NULL, NULL,
            ciphertext, len, ad, ad_len, tag, key, nonce) != 0)
        return -1;  // Authentication failure
    // 2. Only now is plaintext safe to use
    return len;
}
\end{lstlisting}

Never decrypt without verifying the tag. Never release decrypted plaintext before tag verification.

\section{Exercises}

\begin{exercise}
	Given the following constant-time vulnerabilities, identify the leak:
	\begin{lstlisting}[language=C]
int ct_compare(const uint8_t *a, const uint8_t *b, size_t n) {
    volatile uint8_t diff = 0;
    for (size_t i = 0; i < n; i++) {
        diff |= a[i] ^ b[i];
    }
    return diff;
}
	\end{lstlisting}
	Is this truly constant-time? What if \texttt{a} or \texttt{b} is paged out?
\end{exercise}

\begin{exercise}
	Design a constant-time function that selects one of two buffers based on a secret condition: \texttt{void ct\_select(int condition, uint8\_t *out, const uint8\_t *a, const uint8\_t *b, size\_t n)}. If condition is true, copy \texttt{a} to \texttt{out}; otherwise copy \texttt{b}. Both branches must execute regardless of condition.
\end{exercise}

\begin{exercise}
	Why is AES-GCM preferred over AES-CBC + HMAC? What security property does GCM provide that CBC+HMAC requires extra care to achieve?
\end{exercise}

\begin{exercise}
	Research the Debian OpenSSL RNG vulnerability (CVE-2008-0166). What was the root cause? How many cryptographic keys were affected? What lesson does this teach about RNG in systems programming?
\end{exercise}
